% !TEX spellcheck = en_US



% ~~~~~~~~~~~~~~~~~~~~~~~~~~~~~~~~~~~~~~~ V
% (NC) 2026 Haluk Bingol
% github.com/halukbingol/LaTeX-Templates
%
% Licensees may copy, distribute, display, and perform the work and make derivative works 
% and remixes based on it only for non-commercial purposes. 
% ~~~~~~~~~~~~~~~~~~~~~~~~~~~~~~~~~~~~~~~ A




% =====================================================================
%  Generic Stack and Queue -- Junior Level
%  ---------------------------------------------------------------------
%  Built on hblib/Template/hblibTemplate.tex with the libraries in ../../hblib/.
%  Programs and their real outputs are in codeoutput/.
% =====================================================================

\documentclass[aspectratio=169,10pt,compress]{beamer}

% ~~~~~~~~~~~~~~~~~~~~~~~~~~~~~~~~~~~~~~~~~~~~~~~~~~~~~~~~~~~~~~~~~ V hblib
	\usepackage{../../hblib/hblibPresentation}
	\usepackage{../../hblib/hblibCodeoutput}
	\usepackage{../../hblib/hblibDatastructures}
% ~~~~~~~~~~~~~~~~~~~~~~~~~~~~~~~~~~~~~~~~~~~~~~~~~~~~~~~~~~~~~~~~~ A hblib


\usepackage[T1]{fontenc}
\usepackage{lmodern}
\usepackage{booktabs}

\newcommand{\win}{\textbf{Windows}}
\newcommand{\mac}{\textbf{macOS}}

\lstset{columns=fixed}

\SetAlCapFnt{\small}
\SetAlCapNameFnt{\small}
\SetKwProg{Fn}{function}{}{end}
\SetKw{KwError}{error}

% ---------------------------------------------------------------------
\title{Set, Stack and Queue}
\subtitle{Junior Level: Choosing the Right Container}
\author[Bingol]{Ruhan Alpaydin}

\institute[FCIS]{
    Faculty of Computer and Information Sciences\\
    Yeditepe University
}
\date{
	{\tiny \hbTimeStamp}
}

\begin{document}




% ~~~~~~~~~~~~~~~~~~~~~~~~~~~~~~~~~~~~~~ V frm
\begin{frame}[t,fragile]
  \titlepage
\end{frame}
% ~~~~~~~~~~~~~~~~~~~~~~~~~~~~~~~~~~~~~~ A frm




% ~~~~~~~~~~~~~~~~~~~~~~~~~~~~~~~~~~~~~~ V frm
\begin{frame}[t,fragile] \frametitle{Outline}
  \begin{columns}[T]
    \begin{column}{0.48\textwidth}
      \tableofcontents[sections={1-4}, hideallsubsections]
    \end{column}
    \begin{column}{0.48\textwidth}
      \tableofcontents[sections={5-8}, hideallsubsections]
    \end{column}
  \end{columns}
\end{frame}
% ~~~~~~~~~~~~~~~~~~~~~~~~~~~~~~~~~~~~~~ A frm

% ~~~~~~~~~~~~~~~~~~~~~~~~~~~~~~~~~~~~~~ V
\hSection[Ready]{From Building to Choosing}
% ~~~~~~~~~~~~~~~~~~~~~~~~~~~~~~~~~~~~~~ A




% ~~~~~~~~~~~~~~~~~~~~~~~~~~~~~~~~~~~~~~ V frm
\begin{frame}[t,fragile] \frametitle{From building to choosing}
  \begin{itemize}
    \item Freshman: \hCode{new ArrayDeque<>()}, done.
    \item Sophomore: build \hCode{ArrayStack<T>}, \hCode{ArrayQueue<T>}, \hCode{ArraySet<T>}.
    \item Junior: \hRemark{which one}, when, and \emph{why}.
  \end{itemize}
  \begin{block}{This deck answers three questions}
    \small
    \begin{itemize}
      \item \textbf{Array or linked?} Different costs, different caches.
      \item \textbf{Primitive or boxed?} \hCode{int} vs \hCode{Integer}, in a stack or in a queue.
      \item \textbf{Array or heap?} FIFO queue vs \hCode{PriorityQueue}.
    \end{itemize}
  \end{block}
\end{frame}
% ~~~~~~~~~~~~~~~~~~~~~~~~~~~~~~~~~~~~~~ A frm




% ~~~~~~~~~~~~~~~~~~~~~~~~~~~~~~~~~~~~~~ V frm
\begin{frame}[t,fragile] \frametitle{Three backing stores}
  \begin{center}\small
  \begin{tabular}{L{0.20\textwidth}L{0.22\textwidth}L{0.24\textwidth}L{0.24\textwidth}}
    \toprule
    \textbf{Store}  & \textbf{What you pay} & \textbf{What you save} & \textbf{When it wins} \\
    \midrule
    Fixed array     & overflow possible      & nothing to allocate    & small, bounded \\
    Circular array (growing) & copies at $1,2,4,\dots$ & one array, one reference & \textbf{almost always} \\
    Linked list     & one object per element & no copies, no wrap      & splice, iterate, unknown size \\
    \bottomrule
  \end{tabular}
  \end{center}
  \vspace{0.4em}
  \small
  The JDK ships both. \hCode{ArrayDeque} is a growing circular array;
  \hCode{LinkedList} is a doubly linked list. Neither is ``best'':
  each wins on a different workload.
\end{frame}
% ~~~~~~~~~~~~~~~~~~~~~~~~~~~~~~~~~~~~~~ A frm




% ~~~~~~~~~~~~~~~~~~~~~~~~~~~~~~~~~~~~~~ V frm
\begin{frame}[t,fragile] \frametitle{When each wins}
  \begin{columns}[T]
    \begin{column}{0.48\textwidth}
      \begin{block}{Array wins}
        \small
        \begin{itemize}
          \item push / pop / enqueue / dequeue only
          \item \textbf{sequential} iteration
          \item many small elements
          \item the CPU cache can help
        \end{itemize}
      \end{block}
    \end{column}
    \begin{column}{0.48\textwidth}
      \begin{block}{Linked list wins}
        \small
        \begin{itemize}
          \item insert or remove in the \textbf{middle} (iterator)
          \item no reallocation ever
          \item \hCode{null} is allowed (unlike \hCode{ArrayDeque})
        \end{itemize}
      \end{block}
    \end{column}
  \end{columns}
  \vspace{0.6em}
  \small
  For a stack or a queue, the array usually wins.
  The rest of this deck shows \emph{why}.
\end{frame}
% ~~~~~~~~~~~~~~~~~~~~~~~~~~~~~~~~~~~~~~ A frm

% ~~~~~~~~~~~~~~~~~~~~~~~~~~~~~~~~~~~~~~ V
\hSection[Linked]{Linked Backing}
% ~~~~~~~~~~~~~~~~~~~~~~~~~~~~~~~~~~~~~~ A




% ~~~~~~~~~~~~~~~~~~~~~~~~~~~~~~~~~~~~~~ V frm
\begin{frame}[t,fragile] \frametitle{\texttt{LinkedStack<T>}: structure}
  \lstcode[basicstyle=\ttfamily\scriptsize, firstline=3, lastline=10, firstnumber=3]{codeoutput/LinkedStack.java}{Java}
  \begin{itemize}\small
    \item A node has the item and a reference to the \emph{next} node.
    \item \hCode{top} is the most recently pushed node, not an index.
    \item No array, no wrap, no capacity. \hRemark{One} object per element.
    \item Memory cost: the item plus a reference (usually $4$ or $8$ bytes).
  \end{itemize}
\end{frame}
% ~~~~~~~~~~~~~~~~~~~~~~~~~~~~~~~~~~~~~~ A frm




% ~~~~~~~~~~~~~~~~~~~~~~~~~~~~~~~~~~~~~~ V frm
\begin{frame}[t,fragile] \frametitle{\texttt{LinkedStack<T>}: push and pop}
  \lstcode[basicstyle=\ttfamily\scriptsize, firstline=12, lastline=23, firstnumber=12]{codeoutput/LinkedStack.java}{Java}
  \begin{itemize}\small
    \item \hCode{push}: build a new node pointing at the old \hCode{top}, then advance \hCode{top}.
    \item \hCode{pop}: read \hCode{top.item}, advance \hCode{top}, decrement \hCode{size}.
    \item No \emph{loitering} here: \hCode{top} moves and the old node becomes
          unreachable, so the GC handles it.
    \item \hRemark{Constant work in the worst case}: no copy, ever.
  \end{itemize}
\end{frame}
% ~~~~~~~~~~~~~~~~~~~~~~~~~~~~~~~~~~~~~~ A frm




% ~~~~~~~~~~~~~~~~~~~~~~~~~~~~~~~~~~~~~~ V frm
\begin{frame}[t,fragile] \frametitle{\texttt{LinkedQueue<T>}: enqueue and dequeue}
  \lstcode[basicstyle=\ttfamily\scriptsize, firstline=13, lastline=27, firstnumber=13]{codeoutput/LinkedQueue.java}{Java}
  \begin{itemize}\small
    \item \hCode{enqueue}: append at \hCode{tail}; if empty, both \hCode{head} and
          \hCode{tail} become the new node.
    \item \hCode{dequeue}: advance \hCode{head}; if it becomes \hCode{null},
          reset \hCode{tail} too.
    \item Forgetting the \hCode{tail = null} on the last dequeue leaves a
          \emph{dangling} tail: the queue looks empty but the next enqueue
          chains from a stale node.
  \end{itemize}
\end{frame}
% ~~~~~~~~~~~~~~~~~~~~~~~~~~~~~~~~~~~~~~ A frm

% ~~~~~~~~~~~~~~~~~~~~~~~~~~~~~~~~~~~~~~ V
\hSection[Boxing]{Boxing and Bytes}
% ~~~~~~~~~~~~~~~~~~~~~~~~~~~~~~~~~~~~~~ A




% ~~~~~~~~~~~~~~~~~~~~~~~~~~~~~~~~~~~~~~ V frm
\begin{frame}[t,fragile] \frametitle{Boxing: \texttt{int} vs \texttt{Integer}}
  \begin{block}{You cannot put an \texttt{int} in a \texttt{Stack<Integer>}}
    \small
    Generics work only with \textbf{reference types}, so every
    \hCode{push(1)} creates an \hCode{Integer} object (autoboxing)
    and every \hCode{pop()} reads it back (unboxing).
  \end{block}
  \begin{columns}[T]
    \begin{column}{0.48\textwidth}
      \begin{block}{Cheap}
        \small \hCode{-128..127} come from a cached pool:
        no allocation for small values.
      \end{block}
    \end{column}
    \begin{column}{0.48\textwidth}
      \begin{block}{Expensive}
        \small Otherwise, one \hCode{Integer} object per push: more bytes,
        more GC pressure, worse locality.
      \end{block}
    \end{column}
  \end{columns}
\end{frame}
% ~~~~~~~~~~~~~~~~~~~~~~~~~~~~~~~~~~~~~~ A frm




% ~~~~~~~~~~~~~~~~~~~~~~~~~~~~~~~~~~~~~~ V frm
\begin{frame}[t,fragile] \frametitle{Boxing: measured}
  \begin{columns}[T]
    \begin{column}{0.52\textwidth}
      \lstcode[basicstyle=\ttfamily\tiny]{codeoutput/BoxingDemo.java}{Java}
    \end{column}
    \begin{column}{0.44\textwidth}
      \lstterminal[basicstyle=\ttfamily\scriptsize]{codeoutput/BoxingDemo.txt}
      \small
      One million elements, one run on this machine.
      The boxed array is \hRemark{several times} slower to fill and sum.
    \end{column}
  \end{columns}
\end{frame}
% ~~~~~~~~~~~~~~~~~~~~~~~~~~~~~~~~~~~~~~ A frm




% ~~~~~~~~~~~~~~~~~~~~~~~~~~~~~~~~~~~~~~ V frm
\begin{frame}[t,fragile] \frametitle{Bytes per element}
  \begin{columns}[T]
    \begin{column}{0.52\textwidth}
      \lstcode[basicstyle=\ttfamily\tiny]{codeoutput/BytesPerElement.java}{Java}
    \end{column}
    \begin{column}{0.44\textwidth}
      \lstterminal[basicstyle=\ttfamily\scriptsize]{codeoutput/BytesPerElement.txt}
      \small
      Rule of thumb: \hCode{Integer} on a 64-bit JVM is
      about \hRemark{16--20 bytes}, versus $4$ bytes for \hCode{int}.
      That is why \hCode{ArrayList<Integer>} of a million elements
      is much bigger than \hCode{int[]}.
    \end{column}
  \end{columns}
\end{frame}
% ~~~~~~~~~~~~~~~~~~~~~~~~~~~~~~~~~~~~~~ A frm




% ~~~~~~~~~~~~~~~~~~~~~~~~~~~~~~~~~~~~~~ V frm
\begin{frame}[t,fragile] \frametitle{Cache locality}
  \begin{columns}[T]
    \begin{column}{0.52\textwidth}
      \lstcode[basicstyle=\ttfamily\tiny]{codeoutput/CacheLocality.java}{Java}
    \end{column}
    \begin{column}{0.44\textwidth}
      \lstterminal[basicstyle=\ttfamily\scriptsize]{codeoutput/CacheLocality.txt}
      \small
      Arrays keep elements \textbf{next to each other} in memory.
      A linked list hops from one heap object to another:
      the CPU cache is useless because each hop is a different line.
    \end{column}
  \end{columns}
\end{frame}
% ~~~~~~~~~~~~~~~~~~~~~~~~~~~~~~~~~~~~~~ A frm

% ~~~~~~~~~~~~~~~~~~~~~~~~~~~~~~~~~~~~~~ V
\hSection[Bench]{Benchmark and Potential}
% ~~~~~~~~~~~~~~~~~~~~~~~~~~~~~~~~~~~~~~ A




% ~~~~~~~~~~~~~~~~~~~~~~~~~~~~~~~~~~~~~~ V frm
\begin{frame}[t,fragile] \frametitle{Benchmark: \texttt{ArrayDeque} vs \texttt{LinkedStack}}
  \begin{columns}[T]
    \begin{column}{0.52\textwidth}
      \lstcode[basicstyle=\ttfamily\tiny]{codeoutput/BenchStack.java}{Java}
    \end{column}
    \begin{column}{0.44\textwidth}
      \lstterminal[basicstyle=\ttfamily\scriptsize]{codeoutput/BenchStack.txt}
      \small
      Five million pushes and pops, one run on this machine.
      \hRemark{Always measure}: JIT warm-up, GC, and CPU
      change the numbers between runs.
    \end{column}
  \end{columns}
\end{frame}
% ~~~~~~~~~~~~~~~~~~~~~~~~~~~~~~~~~~~~~~ A frm




% ~~~~~~~~~~~~~~~~~~~~~~~~~~~~~~~~~~~~~~ V frm
\begin{frame}[t,fragile] \frametitle{The same numbers, side by side}
  \begin{center}\small
  \begin{tabular}{@{}lrr@{}}
    \toprule
    & \hCode{ArrayDeque<Integer>} & \hCode{LinkedStack<Integer>} \\
    \midrule
    pushes + pops & $5{\cdot}10^6$          & $5{\cdot}10^6$ \\
    wall time (one run) & $\sim 290$\,ms   & $\sim 1480$\,ms \\
    allocations   & grows about $24$ times  & $5{\cdot}10^6$ nodes \\
    GC pressure   & very low                & high \\
    \bottomrule
  \end{tabular}
  \end{center}
  \vspace{0.4em}
  \small
  \emph{One run on this machine.} The ratio $\sim 5\times$ is typical; the
  absolute numbers depend on the JDK, the GC, and the CPU.
  \hRemark{The constant factor of a linked structure is large},
  because each push is an allocation and each pop lets the GC work.
\end{frame}
% ~~~~~~~~~~~~~~~~~~~~~~~~~~~~~~~~~~~~~~ A frm




% ~~~~~~~~~~~~~~~~~~~~~~~~~~~~~~~~~~~~~~ V frm
\begin{frame}[t,fragile] \frametitle{Growth copies: what a push really costs}
  A \hCode{ArrayStack<T>} that starts with capacity $1$ and doubles when full:
  \begin{center}\small
  \begin{tabular}{@{}lll@{}}
    \toprule
    \textbf{Size before push} & \textbf{Cost} & \textbf{Why} \\
    \midrule
    $0, 1, 2, 4, 8, \dots, 2^k$ & copy $2^k$ elements & array is full \\
    anything else               & write one element   & spare slot present \\
    \bottomrule
  \end{tabular}
  \end{center}
  \vspace{0.3em}
  \small
  Over $n$ pushes, total copies:
  \[
    1 + 2 + 4 + \dots + 2^{k} = 2^{k+1} - 1 < 2n .
  \]
  So \hRemark{each push is $O(1)$ amortized}, even though the
  \emph{one} that copies costs $O(n)$.
\end{frame}
% ~~~~~~~~~~~~~~~~~~~~~~~~~~~~~~~~~~~~~~ A frm




% ~~~~~~~~~~~~~~~~~~~~~~~~~~~~~~~~~~~~~~ V frm
\begin{frame}[t,fragile] \frametitle{The potential method, in one page}
  Attach a \textbf{potential} $\Phi$ to the data structure. The
  \emph{amortized} cost of operation $i$ is
  \[
    \hat{c}_i = c_i + \Phi(D_i) - \Phi(D_{i-1}),
  \]
  so the sum of $\hat{c}_i$ over any sequence bounds the real cost.
  \vspace{0.4em}
  \begin{block}{For \texttt{ArrayStack} with doubling}
    \small
    Let $\Phi = 2 \cdot \text{size} - \text{capacity}$ (never negative if
    capacity is at least twice the size, which doubling guarantees).
    \begin{itemize}
      \item Push without copy: $c=1$, $\Delta\Phi=2$, so $\hat{c}=3$.
      \item Push with copy of $k$ elements: $c = 1+k$, $\Delta\Phi = 2 - k$,
            so $\hat{c} = 3$.
    \end{itemize}
    \hRemark{Every push has amortized cost at most 3}, independent of $n$.
  \end{block}
  \vspace{0.2em}
  \small
  A \hCode{LinkedStack} does not need this argument: its push and pop are
  $O(1)$ in the \emph{worst} case, since there is no copy to perform.
\end{frame}
% ~~~~~~~~~~~~~~~~~~~~~~~~~~~~~~~~~~~~~~ A frm

% ~~~~~~~~~~~~~~~~~~~~~~~~~~~~~~~~~~~~~~ V
\hSection[Apps]{The Queue in Real Systems}
% ~~~~~~~~~~~~~~~~~~~~~~~~~~~~~~~~~~~~~~ A




% ~~~~~~~~~~~~~~~~~~~~~~~~~~~~~~~~~~~~~~ V frm
\begin{frame}[t,fragile] \frametitle{Breadth-first search is a queue}
  \begin{columns}[T]
    \begin{column}{0.52\textwidth}
      \lstcode[basicstyle=\ttfamily\tiny]{codeoutput/Bfs.java}{Java}
    \end{column}
    \begin{column}{0.44\textwidth}
      \lstterminal[basicstyle=\ttfamily\scriptsize]{codeoutput/Bfs.txt}
      \small
      Every cell is enqueued once. The queue is the \textbf{frontier}:
      all cells at distance $d$ are processed before any cell at $d+1$.
      The first time the target is dequeued, its distance is the shortest.
    \end{column}
  \end{columns}
\end{frame}
% ~~~~~~~~~~~~~~~~~~~~~~~~~~~~~~~~~~~~~~ A frm




% ~~~~~~~~~~~~~~~~~~~~~~~~~~~~~~~~~~~~~~ V frm
\begin{frame}[t,fragile] \frametitle{Depth-first search is a stack}
  \begin{columns}[T]
    \begin{column}{0.52\textwidth}
      \lstcode[basicstyle=\ttfamily\tiny]{codeoutput/Dfs.java}{Java}
    \end{column}
    \begin{column}{0.44\textwidth}
      \lstterminal[basicstyle=\ttfamily\scriptsize]{codeoutput/Dfs.txt}
      \small
      Same code, \hRemark{same frontier idea}, different container.
      \hCode{push} instead of \hCode{offer}, \hCode{pop} instead of \hCode{poll}.
      The BFS/DFS choice is a container choice.
    \end{column}
  \end{columns}
\end{frame}
% ~~~~~~~~~~~~~~~~~~~~~~~~~~~~~~~~~~~~~~ A frm




% ~~~~~~~~~~~~~~~~~~~~~~~~~~~~~~~~~~~~~~ V frm
\begin{frame}[t,fragile] \frametitle{Producer / consumer: a blocking queue}
  \begin{columns}[T]
    \begin{column}{0.56\textwidth}
      \lstcode[basicstyle=\ttfamily\tiny]{codeoutput/ProducerConsumer.java}{Java}
    \end{column}
    \begin{column}{0.4\textwidth}
      \lstterminal[basicstyle=\ttfamily\scriptsize]{codeoutput/ProducerConsumer.txt}
      \small
      \hCode{ArrayBlockingQueue} (capacity 4): \hCode{put} blocks when full,
      \hCode{take} blocks when empty. Producer and consumer work at different speeds;
      the queue is the buffer between them.
    \end{column}
  \end{columns}
\end{frame}
% ~~~~~~~~~~~~~~~~~~~~~~~~~~~~~~~~~~~~~~ A frm




% ~~~~~~~~~~~~~~~~~~~~~~~~~~~~~~~~~~~~~~ V frm
\begin{frame}[t,fragile] \frametitle{A thread pool is a queue of jobs}
  \begin{columns}[T]
    \begin{column}{0.52\textwidth}
      \lstcode[basicstyle=\ttfamily\tiny]{codeoutput/ThreadPool.java}{Java}
    \end{column}
    \begin{column}{0.44\textwidth}
      \lstterminal[basicstyle=\ttfamily\scriptsize]{codeoutput/ThreadPool.txt}
      \small
      \hCode{newFixedThreadPool(2)} keeps two worker threads waiting on a
      \hCode{BlockingQueue}. Each \hCode{submit} enqueues a job; the first
      free worker dequeues it and runs it.
    \end{column}
  \end{columns}
\end{frame}
% ~~~~~~~~~~~~~~~~~~~~~~~~~~~~~~~~~~~~~~ A frm

% ~~~~~~~~~~~~~~~~~~~~~~~~~~~~~~~~~~~~~~ V
\hSection[Choose]{Choosing the Container}
% ~~~~~~~~~~~~~~~~~~~~~~~~~~~~~~~~~~~~~~ A




% ~~~~~~~~~~~~~~~~~~~~~~~~~~~~~~~~~~~~~~ V frm
\begin{frame}[t,fragile] \frametitle{Bounded vs unbounded}
  \begin{columns}[T]
    \begin{column}{0.5\textwidth}
      \begin{block}{Bounded (fixed capacity)}
        \small
        \begin{itemize}
          \item \hCode{ArrayBlockingQueue}
          \item back-pressure: producer slows down
          \item predictable memory
          \item a full queue is a \emph{signal}, not a bug
        \end{itemize}
      \end{block}
    \end{column}
    \begin{column}{0.46\textwidth}
      \begin{block}{Unbounded (grows)}
        \small
        \begin{itemize}
          \item \hCode{LinkedBlockingQueue}, \hCode{ArrayDeque}
          \item no producer slowdown \dots
          \item \dots but memory can explode under a burst
          \item \hRemark{Hidden DoS}: one producer can starve the heap
        \end{itemize}
      \end{block}
    \end{column}
  \end{columns}
  \vspace{0.5em}
  \small
  Rule of thumb: \hRemark{prefer bounded} on the boundary between two
  subsystems. Unbounded only when you can afford the peak.
\end{frame}
% ~~~~~~~~~~~~~~~~~~~~~~~~~~~~~~~~~~~~~~ A frm




% ~~~~~~~~~~~~~~~~~~~~~~~~~~~~~~~~~~~~~~ V frm
\begin{frame}[t,fragile] \frametitle{\texttt{PriorityQueue} is a heap, not a FIFO queue}
  \begin{columns}[T]
    \begin{column}{0.52\textwidth}
      \lstcode[basicstyle=\ttfamily\tiny]{codeoutput/PriorityQueueDemo.java}{Java}
    \end{column}
    \begin{column}{0.44\textwidth}
      \lstterminal[basicstyle=\ttfamily\scriptsize]{codeoutput/PriorityQueueDemo.txt}
      \small
      Interface \hCode{Queue}, yes. But the element that leaves first is
      the \textbf{smallest}, not the oldest. Real uses:
      Dijkstra, schedulers, top-$k$.
    \end{column}
  \end{columns}
\end{frame}
% ~~~~~~~~~~~~~~~~~~~~~~~~~~~~~~~~~~~~~~ A frm




% ~~~~~~~~~~~~~~~~~~~~~~~~~~~~~~~~~~~~~~ V frm
\begin{frame}[t,fragile] \frametitle{Choosing the container: a decision tree}
  \begin{columns}[T]
    \begin{column}{0.5\textwidth}
      \begin{block}{Order is ``first in, first out''}
        \small
        \begin{itemize}
          \item bounded? \hCode{ArrayBlockingQueue}
          \item else \hCode{ArrayDeque}
          \item splice in the middle? \hCode{LinkedList}
        \end{itemize}
      \end{block}
      \begin{block}{Order is ``last in, first out''}
        \small
        \hCode{ArrayDeque} with \hCode{push/pop}.
      \end{block}
    \end{column}
    \begin{column}{0.46\textwidth}
      \begin{block}{Order is ``priority''}
        \small
        \hCode{PriorityQueue}, or \hCode{PriorityBlockingQueue} for threads.
      \end{block}
      \begin{block}{Order is ``no duplicates''}
        \small
        \begin{itemize}
          \item unordered: \hCode{HashSet}
          \item insertion order: \hCode{LinkedHashSet}
          \item sorted: \hCode{TreeSet}
        \end{itemize}
      \end{block}
    \end{column}
  \end{columns}
\end{frame}
% ~~~~~~~~~~~~~~~~~~~~~~~~~~~~~~~~~~~~~~ A frm




% ~~~~~~~~~~~~~~~~~~~~~~~~~~~~~~~~~~~~~~ V frm
\begin{frame}[t,fragile] \frametitle{Which Java class, in one table}
  \begin{center}\footnotesize
  \begin{tabular}{@{}lL{0.18\textwidth}L{0.24\textwidth}L{0.28\textwidth}@{}}
    \toprule
    \textbf{Interface} & \textbf{Default} & \textbf{Alternatives} & \textbf{Pick it when} \\
    \midrule
    \hCode{Set<E>}   & \hCode{HashSet}        & \hCode{LinkedHashSet}, \hCode{TreeSet} & membership tests \\
    \hCode{Deque<E>} & \hCode{ArrayDeque}     & \hCode{LinkedList}                    & stack or FIFO queue \\
    \hCode{Queue<E>} & \hCode{ArrayDeque}     & \hCode{LinkedList}, \hCode{ArrayBlockingQueue} & unbounded or bounded queue \\
    \hCode{BlockingQueue<E>} & \hCode{LinkedBlockingQueue} & \hCode{ArrayBlockingQueue} & threads on both ends \\
    \hCode{Queue<E>} (priority) & \hCode{PriorityQueue} & \hCode{PriorityBlockingQueue} & smallest (or comparator) first \\
    \bottomrule
  \end{tabular}
  \end{center}
  \vspace{0.3em}
  \small
  \hRemark{Reach for the interface} (\hCode{Set}, \hCode{Queue}, \hCode{Deque}),
  not the class, so you can swap one for another without editing the
  caller.
\end{frame}
% ~~~~~~~~~~~~~~~~~~~~~~~~~~~~~~~~~~~~~~ A frm

% ~~~~~~~~~~~~~~~~~~~~~~~~~~~~~~~~~~~~~~ V
\hSection[Pitfalls]{Three Pitfalls You Will Meet}
% ~~~~~~~~~~~~~~~~~~~~~~~~~~~~~~~~~~~~~~ A




% ~~~~~~~~~~~~~~~~~~~~~~~~~~~~~~~~~~~~~~ V frm
\begin{frame}[t,fragile] \frametitle{Pitfall 1: boxing NPE}
  \begin{columns}[T]
    \begin{column}{0.52\textwidth}
      \lstcode[basicstyle=\ttfamily\tiny]{codeoutput/BoxingNPE.java}{Java}
    \end{column}
    \begin{column}{0.44\textwidth}
      \lstterminal[basicstyle=\ttfamily\scriptsize]{codeoutput/BoxingNPE.txt}
      \small
      \hCode{a} is \hCode{null}; \hCode{a + 1} unboxes it and throws.
      Autoboxing hides the reference, so the NPE is not obvious at the call.
    \end{column}
  \end{columns}
\end{frame}
% ~~~~~~~~~~~~~~~~~~~~~~~~~~~~~~~~~~~~~~ A frm




% ~~~~~~~~~~~~~~~~~~~~~~~~~~~~~~~~~~~~~~ V frm
\begin{frame}[t,fragile] \frametitle{Pitfall 2: mutable elements in a Set}
  \begin{columns}[T]
    \begin{column}{0.52\textwidth}
      \lstcode[basicstyle=\ttfamily\tiny]{codeoutput/NullInSet.java}{Java}
    \end{column}
    \begin{column}{0.44\textwidth}
      \lstterminal[basicstyle=\ttfamily\scriptsize]{codeoutput/NullInSet.txt}
      \small
      \hCode{HashSet} stores objects by \hCode{hashCode}; a mutable element
      whose hash changes after insertion is \emph{lost}.
      Use immutable elements as keys.
    \end{column}
  \end{columns}
\end{frame}
% ~~~~~~~~~~~~~~~~~~~~~~~~~~~~~~~~~~~~~~ A frm




% ~~~~~~~~~~~~~~~~~~~~~~~~~~~~~~~~~~~~~~ V frm
\begin{frame}[t,fragile] \frametitle{Pitfall 3: tests and measurements}
  \begin{columns}[T]
    \begin{column}{0.48\textwidth}
      \begin{block}{JUnit: \emph{Not run here}}
        \small
        \begin{lstlisting}[language=Java, basicstyle=\ttfamily\tiny, frame=single]
@Test
void popReducesSize() {
    ArrayStack<Integer> s = new ArrayStack<>(4);
    s.push(1); s.push(2);
    assertEquals(2, s.size());
    assertEquals(2, s.pop());
    assertEquals(1, s.size());
}
        \end{lstlisting}
      \end{block}
    \end{column}
    \begin{column}{0.48\textwidth}
      \begin{block}{JMH: \emph{Not run here}}
        \small
        \begin{lstlisting}[language=Java, basicstyle=\ttfamily\tiny, frame=single]
@Benchmark
public int arrayDequePushPop() {
    Deque<Integer> s = new ArrayDeque<>();
    for (int i = 0; i < 1000; i++) s.push(i);
    int sum = 0;
    while (!s.isEmpty()) sum += s.pop();
    return sum;
}
        \end{lstlisting}
      \end{block}
    \end{column}
  \end{columns}
  \vspace{0.3em}
  \small
  \hRemark{Both tools fight warm-up}. JUnit warms up the JIT by doing
  many small runs; JMH is designed for benchmarks that are
  \emph{repeatable}. Use JUnit for \emph{correctness},
  JMH for \emph{speed}.
\end{frame}
% ~~~~~~~~~~~~~~~~~~~~~~~~~~~~~~~~~~~~~~ A frm

% ~~~~~~~~~~~~~~~~~~~~~~~~~~~~~~~~~~~~~~ V
\hSection[Wrap]{Wrap-up}
% ~~~~~~~~~~~~~~~~~~~~~~~~~~~~~~~~~~~~~~ A




% ~~~~~~~~~~~~~~~~~~~~~~~~~~~~~~~~~~~~~~ V frm
\begin{frame}[t,fragile] \frametitle{Big ideas}
  \begin{itemize}
    \item For a stack or a queue, the \textbf{circular array} is
          usually faster, smaller, and more cache friendly than a
          linked list.
    \item A \hCode{LinkedStack} is $O(1)$ in the worst case; an
          \hCode{ArrayStack} is $O(1)$ \emph{amortized}, and every
          growth copies the live elements.
    \item The potential method is the formal justification of ``amortized''.
    \item Boxing costs memory ($\sim 16$--$20$ bytes per \hCode{Integer})
          and cache locality; keep primitive arrays for hot loops.
    \item \hRemark{BFS uses a queue, DFS uses a stack}: the same
          frontier idea, two containers.
    \item Producer/consumer and thread pools are \hCode{BlockingQueue}
          workloads; prefer \textbf{bounded} on a boundary.
    \item \hCode{PriorityQueue} is a \hCode{Queue} by interface, but the
          order is by key, not by arrival.
    \item Reach for the \emph{interface}, not the class.
  \end{itemize}
\end{frame}
% ~~~~~~~~~~~~~~~~~~~~~~~~~~~~~~~~~~~~~~ A frm




% ~~~~~~~~~~~~~~~~~~~~~~~~~~~~~~~~~~~~~~ V frm
\begin{frame}[t,fragile] \frametitle{Check yourself}
  \begin{enumerate}\small
    \item Why is \hCode{push} on an \hCode{ArrayStack} amortized
          $O(1)$ and not worst-case $O(1)$?
    \item Where does the potential $\Phi = 2 \cdot \text{size} - \text{capacity}$
          come from, and why is it never negative?
    \item Why is an \hCode{Integer} about $16$--$20$ bytes and an
          \hCode{int} exactly $4$?
    \item Why is a linked traversal slower than an array traversal
          even when both are $O(n)$?
    \item When would you choose \hCode{LinkedList} over \hCode{ArrayDeque}?
    \item Why is \hCode{PriorityQueue} not FIFO?
    \item What is the difference between a bounded and an unbounded queue,
          in terms of failure?
  \end{enumerate}
\end{frame}
% ~~~~~~~~~~~~~~~~~~~~~~~~~~~~~~~~~~~~~~ A frm




% ~~~~~~~~~~~~~~~~~~~~~~~~~~~~~~~~~~~~~~ V frm
\begin{frame}[t,fragile]
	\begin{center}
		\vspace{4cm}
		\hRemark{\Huge Questions?}
	\end{center}
\end{frame}
% ~~~~~~~~~~~~~~~~~~~~~~~~~~~~~~~~~~~~~~ A frm




\end{document}